% =====================================================================
%  Промежуточный отчёт. Этап 3 ВКР
%  Программная реализация прототипа
%
%  Сборка: xelatex -> biber -> xelatex -> xelatex
%  Преамбула и библиография переиспользуются из ../Stage1
% =====================================================================
\documentclass[14pt, a4paper]{extarticle}

% =====================================================================
%  Подключаемые пакеты LaTeX
%  Сборка: XeLaTeX + Biber
%  Соответствие: ГОСТ 7.32-2017, ГОСТ Р 7.0.100-2018
% =====================================================================

% Кодировка и язык
\usepackage{fontspec}
\setmainfont{Times New Roman}
\setsansfont{Arial}
\setmonofont{Courier New}

\usepackage{polyglossia}
\setdefaultlanguage{russian}
\setotherlanguage{english}

% Геометрия страницы (ГОСТ 7.32-2017, п. 6.1):
% левое 30 мм, правое 15 мм, верхнее 20 мм, нижнее 20 мм
\usepackage[
    a4paper,
    left=30mm,
    right=15mm,
    top=20mm,
    bottom=20mm,
    footskip=10mm
]{geometry}

% Межстрочный интервал 1.5 (ГОСТ 7.32-2017, п. 6.3)
\usepackage{setspace}
\onehalfspacing

% Абзацный отступ 1.25 см
\usepackage{indentfirst}
\setlength{\parindent}{1.25cm}

% Математика
\usepackage{amsmath}
\usepackage{amssymb}
\usepackage{amsthm}
\usepackage{mathtools}

% Графика и таблицы
\usepackage{graphicx}
\usepackage{booktabs}
\usepackage{longtable}
\usepackage{tabularx}
\usepackage{array}
\usepackage{multirow}
\usepackage{makecell}
\usepackage{caption}
\usepackage{subcaption}
\usepackage{float}

% Цвет, ссылки, листинги
\usepackage{xcolor}
\usepackage[hidelinks]{hyperref}
\hypersetup{
    pdfencoding=auto,
    bookmarksnumbered=true,
    bookmarksopen=true,
    colorlinks=false
}
\usepackage{listings}
\lstset{
    basicstyle=\ttfamily\footnotesize,
    breaklines=true,
    showstringspaces=false,
    keywordstyle=\bfseries,
    columns=fullflexible,
    captionpos=b,
    frame=single,
    aboveskip=8pt,
    belowskip=8pt
}

% Перечисления и списки
\usepackage{enumitem}
\setlist{nosep, leftmargin=*}

% Заголовки и оглавление
\usepackage{titlesec}
\usepackage{tocloft}

% Библиография по ГОСТ Р 7.0.100-2018
\usepackage[
    backend=biber,
    style=gost-numeric,
    language=auto,
    sorting=none,
    maxbibnames=4,
    minbibnames=1
]{biblatex}

% Дополнительные служебные пакеты
\usepackage{csquotes}
\usepackage{etoolbox}
\usepackage{url}
\usepackage{microtype}
\usepackage{xurl}

% =====================================================================
%  Оформление по ГОСТ 7.32-2017
%  Структурные элементы: разделы, рисунки, таблицы, формулы, списки
% =====================================================================

% Размер основного шрифта 14 pt с межстрочным интервалом 1.5 устанавливается
% в classfile (extarticle/extreport, базовый размер 14pt). Здесь — стилевые
% настройки заголовков и нумерации.

% Заголовки по ГОСТ 7.32-2017 (п. 6.2):
%   - заголовки разделов и подразделов с абзацного отступа,
%     с прописной буквы, без точки в конце, не подчёркиваются;
%   - расстояние между заголовком и текстом — не менее одного интервала;
%   - заголовки разделов набираются жирным шрифтом, размер 14 pt.

\titleformat{\section}
    {\normalfont\fontsize{14}{17}\bfseries}
    {\thesection}{1em}{}
\titlespacing*{\section}{\parindent}{18pt}{12pt}

\titleformat{\subsection}
    {\normalfont\fontsize{14}{17}\bfseries}
    {\thesubsection}{1em}{}
\titlespacing*{\subsection}{\parindent}{14pt}{10pt}

\titleformat{\subsubsection}
    {\normalfont\fontsize{14}{17}\bfseries}
    {\thesubsubsection}{1em}{}
\titlespacing*{\subsubsection}{\parindent}{12pt}{8pt}

% Нумерация рисунков и таблиц — сквозная в пределах раздела
\renewcommand{\thefigure}{\arabic{section}.\arabic{figure}}
\renewcommand{\thetable}{\arabic{section}.\arabic{table}}
\renewcommand{\theequation}{\arabic{section}.\arabic{equation}}

% Подписи рисунков и таблиц по ГОСТ:
%   - «Рисунок N — Название» (тире), снизу, по центру;
%   - «Таблица N — Название», сверху, слева.
\captionsetup[figure]{
    labelsep=endash,
    name=Рисунок,
    justification=centering,
    singlelinecheck=true,
    font=normalsize
}
\captionsetup[table]{
    labelsep=endash,
    name=Таблица,
    justification=raggedright,
    singlelinecheck=false,
    font=normalsize
}

% Колонтитулы — простые, номер страницы по центру нижнего поля
\usepackage{fancyhdr}
\pagestyle{fancy}
\fancyhf{}
\fancyfoot[C]{\thepage}
\renewcommand{\headrulewidth}{0pt}
\renewcommand{\footrulewidth}{0pt}

% Оглавление: «СОДЕРЖАНИЕ» по центру, прописными
\renewcommand{\contentsname}{\hfill СОДЕРЖАНИЕ\hfill}

% Список литературы: «СПИСОК ИСПОЛЬЗОВАННЫХ ИСТОЧНИКОВ»
\defbibheading{bibliography}[Список использованных источников]{%
    \section*{#1}\addcontentsline{toc}{section}{#1}%
}

% =====================================================================
%  Пользовательские команды и сокращения
% =====================================================================

% Краткие ссылки
\newcommand{\figref}[1]{рисунок~\ref{#1}}
\newcommand{\Figref}[1]{Рисунок~\ref{#1}}
\newcommand{\tabref}[1]{таблица~\ref{#1}}
\newcommand{\Tabref}[1]{Таблица~\ref{#1}}
\newcommand{\eqnref}[1]{(\ref{#1})}
\newcommand{\secref}[1]{раздел~\ref{#1}}

% Английские сокращения как термины
\newcommand{\eng}[1]{\textit{#1}}

% Метки общих обозначений
\DeclareMathOperator*{\argmax}{arg\,max}
\DeclareMathOperator*{\argmin}{arg\,min}
\DeclareMathOperator{\softmax}{softmax}
\DeclareMathOperator{\sigm}{\sigma}
\DeclareMathOperator{\KL}{KL}
\DeclareMathOperator{\MMD}{MMD}
\DeclareMathOperator{\PSI}{PSI}

% Векторы и матрицы — обычным начертанием в соответствии с ГОСТ 7.32
% (без полужирного, как принято в большинстве отечественных НИР)

% Окружение для алгоритмов
\newtheoremstyle{algodef}{8pt}{8pt}{\normalfont}{}{\bfseries}{.}{0.5em}{}
\theoremstyle{algodef}
\newtheorem{algorithm}{Алгоритм}[section]
\newtheorem{definition}{Определение}[section]
\newtheorem{statement}{Утверждение}[section]


\addbibresource{../Stage1/bib/refs.bib}

\hypersetup{
    pdftitle={Промежуточный отчёт. Этап 3 ВКР},
    pdfsubject={Программная реализация прототипа NIDS},
    pdfkeywords={CNN-LSTM, PyTorch, FastAPI, Streamlit, UNSW-NB15,
                 ИИ-злоумышленник, drift, FSM-агент}
}

\begin{document}

% ----------- Титульный лист -----------
\begin{titlepage}
\centering
\vspace*{1cm}
{\large\bfseries Промежуточный отчёт по выпускной квалификационной работе}\\[0.3cm]
{\large\bfseries Этап 3. Программная реализация прототипа}\\[2cm]
{\Large\bfseries Тема: <<Методика обнаружения несанкционированных действий
и аномальной активности в корпоративной сети с использованием нейронных
сетей>>}\\[3cm]
\begin{flushright}
\begin{tabular}{rl}
Исполнитель: & \rule{6cm}{0.4pt} \\[0.3cm]
Научный руководитель: & \rule{6cm}{0.4pt} \\
\end{tabular}
\end{flushright}
\vfill
{\large \the\year}\\
\end{titlepage}

% ----------- Реферат -----------
\section*{РЕФЕРАТ}
\addcontentsline{toc}{section}{РЕФЕРАТ}

\textbf{Ключевые слова:} прототип, Python, PyTorch, scikit-learn, XGBoost,
FastAPI, Streamlit, CNN-LSTM, UNSW-NB15, CICIDS2017, ИИ-злоумышленник,
конечный автомат, drift-инжектор, sliding window, focal loss, температурное
масштабирование, PSI, MMD.

\textbf{Цель этапа.} Программная реализация всех компонентов системы,
зафиксированных в Этапах~1--2, формирование прототипа, готового к
проведению экспериментов Этапа~4, и описание архитектуры реализации в
формате, позволяющем воспроизвести эксперименты на стороннем стенде.

\textbf{Содержание этапа.} В рамках этапа реализованы: подсистема
предобработки UNSW-NB15 с воспроизводимым трансформером \texttt{Preprocessor};
14 моделей-кандидатов (CNN-LSTM как основная, плюс альтернативы и baselines);
training/evaluation pipeline с focal-loss, температурным масштабированием
и bootstrap-CI; ИИ-злоумышленник в виде 7 параметризованных шаблонов,
конечного автомата с одиннадцатью состояниями и drift-инжектора трёх типов
(covariate, prior, concept); realtime-стенд на FastAPI и Streamlit. Покрытие
исходного кода юнит-тестами включает не менее 7 тестовых модулей.

\newpage

% ----------- Содержание -----------
\tableofcontents
\newpage

\section*{ВВЕДЕНИЕ}
\addcontentsline{toc}{section}{ВВЕДЕНИЕ}

В настоящем отчёте описана программная реализация прототипа системы
обнаружения несанкционированных действий, выполненная на основании
проектных решений Этапов~1 и 2. Все архитектурные решения, обоснованные
ранее (CNN-LSTM как proposed; UNSW-NB15 как основной датасет;
шаблоны+FSM+drift как ИИ-злоумышленник без cGAN; FastAPI+Streamlit как
интерфейс), реализованы и покрыты тестами. Реализация ориентирована на
воспроизводимость, контролируется единым набором конфигурационных файлов
и собирается одной командой Make.

\section{Структура программного проекта}

Проект организован как Python-пакет \texttt{diploma\_nids} с
исчерпывающим разделением ответственности по подпакетам:
\texttt{data}, \texttt{models}, \texttt{training}, \texttt{eval},
\texttt{attacker}, \texttt{inference}, \texttt{ui}, \texttt{utils}.
Сводное представление структуры приведено в \tabref{tab:project-layout}.

\begin{table}[H]
\centering
\caption{Структура программного проекта}
\label{tab:project-layout}
\renewcommand{\arraystretch}{1.2}
\setlength{\tabcolsep}{4pt}
\small
\begin{tabularx}{\linewidth}{@{}lXX@{}}
\toprule
\textbf{Каталог} & \textbf{Назначение} & \textbf{Состав (ключевые модули)} \\
\midrule
\texttt{configs/}      & Параметризация всех стадий pipeline &
data, preprocess, models, train, eval, attacker, pipeline (15 YAML) \\
\addlinespace
\texttt{src/diploma\_nids/data/} & Схемы, загрузка, препроцессинг, окна, splits &
\texttt{schema.py}, \texttt{load.py}, \texttt{preprocess.py},
\texttt{windowing.py}, \texttt{splits.py} \\
\addlinespace
\texttt{src/diploma\_nids/models/} & Модели-кандидаты &
\texttt{base.py} (registry), \texttt{cnn\_lstm.py},
\texttt{rnn\_family.py}, \texttt{cnn1d.py}, \texttt{tcn.py},
\texttt{transformer.py}, \texttt{autoencoder.py}, \texttt{mlp.py},
\texttt{classical.py} \\
\addlinespace
\texttt{src/diploma\_nids/training/} & Обучение DL-моделей &
\texttt{trainer.py}, \texttt{losses.py} (BCE, focal) \\
\addlinespace
\texttt{src/diploma\_nids/eval/} & Метрики, калибровка, drift &
\texttt{metrics.py}, \texttt{thresholding.py}, \texttt{drift.py},
\texttt{error\_analysis.py} \\
\addlinespace
\texttt{src/diploma\_nids/attacker/} & ИИ-злоумышленник &
\texttt{templates/} (7 шаблонов + normal), \texttt{agent.py} (FSM),
\texttt{drift\_injector.py}, \texttt{runtime.py} \\
\addlinespace
\texttt{src/diploma\_nids/inference/} & Online-инференс и алертинг &
\texttt{stream.py} (WindowScorer), \texttt{alerting.py} (AlertFormer),
\texttt{service.py} (FastAPI) \\
\addlinespace
\texttt{src/diploma\_nids/ui/} & Streamlit-дэшборд &
\texttt{streamlit\_app.py} \\
\addlinespace
\texttt{src/diploma\_nids/utils/} & Воспроизводимость и IO &
\texttt{seed.py}, \texttt{io.py}, \texttt{logging.py} \\
\addlinespace
\texttt{scripts/} & Точки входа &
\texttt{01\_download\_unsw.py}\dots\texttt{07\_demo\_realtime.py} \\
\addlinespace
\texttt{tests/} & Юнит-тесты &
\texttt{test\_utils.py}, \texttt{test\_configs.py},
\texttt{test\_preprocess.py}, \texttt{test\_windowing.py},
\texttt{test\_splits.py}, \texttt{test\_models\_smoke.py},
\texttt{test\_metrics.py}, \texttt{test\_attacker.py},
\texttt{test\_alerting.py} \\
\bottomrule
\end{tabularx}
\end{table}

Сборка и операционный pipeline инкапсулированы в \texttt{Makefile}: цели
\texttt{install}, \texttt{install-dev}, \texttt{lint}, \texttt{type},
\texttt{test}, \texttt{data}, \texttt{preprocess}, \texttt{train},
\texttt{evaluate}, \texttt{attack}, \texttt{drift-test}, \texttt{demo},
\texttt{reproduce}.

\section{Подсистема обработки данных}

\subsection{Схема и загрузка}

Схема UNSW-NB15 формализована pydantic-моделями (модуль
\texttt{data/schema.py}) и описывает 5 групп признаков, согласованных с
\tabref{tab:unsw-feature-groups} Этапа~2. Загрузчик
(\texttt{data/load.py}) валидирует наличие признаков и меток, проверяет
существование официального train/test-разделения и возвращает
\texttt{DataFrame}. Аналогично реализована загрузка CICIDS2017 для
cross-evaluation.

\subsection{Препроцессор}

Класс \texttt{Preprocessor} реализует sklearn-совместимый интерфейс
(\texttt{fit}/\texttt{transform}/\texttt{fit\_transform}). На стадии
\texttt{fit} (на обучающей выборке) формируются:

\begin{itemize}
    \item квантильные границы для обрезки выбросов
    (\(0{,}001\) и \(0{,}999\) по умолчанию);
    \item статистики robust-scaler (медиана, IQR) для каждого числового
    признака;
    \item one-hot категории для признаков с малой кардинальностью
    (\(\leq 16\)) и frequency- или target-encoding для большой кардинальности
    (\emph{service});
    \item фиксированный порядок выходных столбцов \texttt{output\_columns},
    обеспечивающий идентичность представления в train/val/test/inference.
\end{itemize}

Все стадии описываются формулами Stage~2 §2.4--2.7. Препроцессор
сериализуется в JSON (метод \texttt{save}) и восстанавливается
(\texttt{load}) для онлайн-инференса.

\subsection{Скользящие окна}

Модуль \texttt{data/windowing.py} реализует \texttt{build\_windows}, который
по векторизованному массиву признаков и меткам возвращает тензор
\((m, W, d)\) и метки окон по одной из стратегий агрегации (\emph{any},
\emph{majority}, \emph{last}). Параметры по умолчанию: \(W = 32\),
\(S = W/4 = 8\). Запрет утечки достигается тем, что окна формируются
независимо для каждого split-а после временного разделения.

\section{Подсистема моделей}

\subsection{Регистрация и единый интерфейс}

Базовый модуль \texttt{models/base.py} вводит \texttt{BaseDetector} (для
классических моделей с интерфейсом \texttt{fit}/\texttt{predict\_proba}/
\texttt{score}/\texttt{predict}), \texttt{DLDetector} (mix-in для torch-моделей
с унифицированными \texttt{score}/\texttt{save}/\texttt{load}) и реестр
\texttt{register}/\texttt{get\_model}/\texttt{build\_model}. Регистрация
выполняется при импорте подпакета \texttt{models}.

\subsection{Реализованные модели}

В \tabref{tab:models-impl} сведены реализованные модели с указанием
файлов и ключевых деталей.

\begin{table}[H]
\centering
\caption{Реализованные модели}
\label{tab:models-impl}
\renewcommand{\arraystretch}{1.2}
\small
\begin{tabularx}{\linewidth}{@{}llXl@{}}
\toprule
\textbf{Имя} & \textbf{Тип} & \textbf{Файл} & \textbf{Окна} \\
\midrule
mlp                  & DL supervised        & \texttt{models/mlp.py}        & нет \\
cnn1d                & DL supervised        & \texttt{models/cnn1d.py}      & да  \\
\textbf{cnn\_lstm}   & DL supervised \emph{(proposed)} & \texttt{models/cnn\_lstm.py} & да \\
lstm                 & DL supervised        & \texttt{models/rnn\_family.py} & да \\
gru                  & DL supervised        & \texttt{models/rnn\_family.py} & да \\
bilstm               & DL supervised + attn & \texttt{models/rnn\_family.py} & да \\
tcn                  & DL supervised        & \texttt{models/tcn.py}        & да  \\
transformer          & DL supervised        & \texttt{models/transformer.py} & да \\
autoencoder          & DL semi-supervised   & \texttt{models/autoencoder.py} & нет \\
vae                  & DL semi-supervised   & \texttt{models/autoencoder.py} & нет \\
logistic\_regression & classical            & \texttt{models/classical.py}  & нет \\
random\_forest       & classical            & \texttt{models/classical.py}  & нет \\
xgboost              & classical            & \texttt{models/classical.py}  & нет \\
isolation\_forest    & classical (unsup.)   & \texttt{models/classical.py}  & нет \\
ocsvm                & classical (one-class)& \texttt{models/classical.py}  & нет \\
\bottomrule
\end{tabularx}
\end{table}

Все DL-модели реализованы на PyTorch и единообразно возвращают объект
\texttt{ModelOutput} с полями \texttt{logits}, \texttt{probs} и
\texttt{extras} для дополнительных тензоров (например, реконструкции и
KL для VAE). Архитектура CNN-LSTM строго следует \secref{subsec:why-cnnlstm}
Этапа~2: блоки Conv1D-BN-ReLU-Pool-Dropout, LSTM, MLP-голова, sigmoid-логит.

\section{Подсистема обучения и оценки}

\subsection{Loss-функции}

В модуле \texttt{training/losses.py} реализованы binary cross-entropy с
логитами и фокальная функция потерь~\eqnref{eq:focal} (Lin и
соавт.~\cite{lin2017focal}). Параметры \(\alpha\), \(\gamma\) задаются в
\texttt{configs/train/default.yaml} (значения по умолчанию
\(\alpha = 0{,}25\), \(\gamma = 2{,}0\)).

\subsection{Тренер}

\texttt{training/trainer.py} реализует:
\begin{itemize}
    \item AdamW + cosine-LR с min\_lr (опционально step-LR);
    \item gradient clipping;
    \item early stopping по \texttt{val\_pr\_auc} (по умолчанию max-mode,
    patience = 10);
    \item чекпойнт лучшей модели и журнал history (train/val loss,
    val\_pr\_auc, val\_f1) на каждом эпохе.
\end{itemize}

\subsection{Метрики и калибровка}

\texttt{eval/metrics.py} реализует Precision, Recall, F1, FPR, MCC,
ROC-AUC, PR-AUC, ECE, FP-per-time, time-to-detect, поиск порога по
целевому FPR и по F1-max, bootstrap-CI. Модуль
\texttt{eval/thresholding.py} содержит температурное
масштабирование~\cite{guo2017calibration} с обучением через L-BFGS на
валидационных логитах.

\subsection{Drift-метрики}

\texttt{eval/drift.py} реализует PSI~\eqnref{eq:psi}, KL и MMD с
автобандвидтом по медианной эвристике. Функция \texttt{drift\_report}
возвращает агрегат с явным индикатором \texttt{drift\_alarm} при
превышении PSI-порога \(0{,}25\), используется в Эксперименте~5.

\section{ИИ-злоумышленник}

\subsection{Шаблоны атак}

В \texttt{attacker/templates/} реализованы 7 параметризованных шаблонов:
\texttt{ddos}, \texttt{scan}, \texttt{brute}, \texttt{exploit}, \texttt{exfil},
\texttt{intrec}, \texttt{worm}, плюс \texttt{normal}. Каждый шаблон строит
последовательность объектов \texttt{FlowRecord} (датакласс, эквивалентный
49 признакам UNSW-NB15) с явной привязкой к классу \emph{attack\_cat} из
\tabref{tab:attack-templates} Этапа~2.

\subsection{Конечно-автоматный планировщик}

\texttt{attacker/agent.py} реализует \texttt{FSMAgent} с поддержкой
одиннадцати состояний (\emph{NORMAL}, 7 \emph{ATTACK} и 3 \emph{DRIFT}),
вероятностями переходов из \texttt{configs/attacker/policy.yaml} и
минимальными длительностями состояний. Каждый тик возвращает структуру
\texttt{AgentTick} с метаданными ground truth (тип атаки, состояние FSM,
тип дрейфа), необходимыми для расчёта метрик в Этапе~4.

\subsection{Drift-инжектор}

\texttt{attacker/drift\_injector.py} реализует три типа дрейфа:
\begin{itemize}
    \item \emph{covariate} --- сдвиг масштаба и шума по выбранным
    числовым признакам с коэффициентами из конфига;
    \item \emph{prior} --- управляется через политику FSM (распределение
    долей \emph{ATTACK}-состояний);
    \item \emph{concept} --- маскирование сигнатурных признаков атаки
    (например, \texttt{sload}, \texttt{dload}) с интенсивностью,
    нарастающей по \emph{ramp\_up\_seconds}.
\end{itemize}

\subsection{Runtime}

\texttt{attacker/runtime.py} (\texttt{AttackerRuntime}) объединяет агент
и drift-инжектор и предоставляет два режима работы: offline (метод
\texttt{collect\_dataframe}) для Эксперимента~4 и online (метод
\texttt{stream\_with\_callback}) для realtime-демо.

\section{Подсистема online-инференса и интерфейс}

\subsection{FastAPI-сервис}

Модуль \texttt{inference/service.py} реализует REST API со следующими
конечными точками:
\begin{itemize}
    \item \texttt{GET /health} --- проверка живости;
    \item \texttt{GET /info} --- сведения о загруженной модели;
    \item \texttt{POST /score} --- скоринг батча flow-записей; для
    моделей с \texttt{expects\_windows = True} используется
    скользящий буфер \texttt{WindowScorer};
    \item \texttt{POST /score-window} --- прямой скоринг готового окна
    \((W, F)\);
    \item \texttt{GET /alerts/recent} --- последние \(n\) алертов.
\end{itemize}

Загрузка артефактов происходит из переменных окружения
\texttt{DIPLOMA\_MODEL\_YAML}, \texttt{DIPLOMA\_CHECKPOINT},
\texttt{DIPLOMA\_PREPROCESSOR}, \texttt{DIPLOMA\_THRESHOLD}, что
обеспечивает декларативность конфигурации без изменения кода.

\subsection{Streamlit-дэшборд}

\texttt{ui/streamlit\_app.py} реализует обновляемую с интервалом
1--2~с панель: метрики (число алертов, доля \emph{critical/high},
средний скоринг), таблицу алертов и временной ряд скорингов. Сообщается с
сервисом по \texttt{DIPLOMA\_API\_URL}.

\subsection{Demo-orchestrator}

\texttt{scripts/07\_demo\_realtime.py} запускает FastAPI и Streamlit как
дочерние процессы и стримит выход \texttt{AttackerRuntime} в
\texttt{POST /score} в реальном времени. Это и есть стенд Эксперимента~4
из протокола Этапа~2.

\section{Воспроизводимость и журналирование}

\begin{itemize}
    \item Seed фиксируется в каждой точке входа через
    \texttt{utils/seed.py} (\texttt{numpy}, \texttt{random},
    \texttt{torch}, \texttt{torch.cuda}, \texttt{PYTHONHASHSEED}, и
    \texttt{torch.use\_deterministic\_algorithms}).
    \item Все артефакты обучения и оценки сохраняются в виде JSON-файлов
    в \texttt{experiments/runs/}, чекпойнты --- в \texttt{models/}.
    \item Конфигурации экспериментов фиксируются YAML-ями в
    \texttt{configs/}; \texttt{Makefile reproduce} запускает
    \texttt{data $\to$ preprocess $\to$ train $\to$ evaluate} одной
    командой.
    \item Препроцессор сохраняется JSON-ом и применяется идентично в
    обучении и инференсе, исключая test-time leakage.
\end{itemize}

\section{Покрытие тестами}

В \tabref{tab:tests} приведён состав юнит-тестов и проверяемые
свойства.

\begin{table}[H]
\centering
\caption{Юнит-тесты прототипа}
\label{tab:tests}
\renewcommand{\arraystretch}{1.2}
\small
\begin{tabularx}{\linewidth}{@{}lX@{}}
\toprule
\textbf{Файл} & \textbf{Проверяемые свойства} \\
\midrule
\texttt{test\_utils.py}    & seed-воспроизводимость; YAML/JSON roundtrip; создание родительских директорий \\
\texttt{test\_configs.py}  & все YAML-ы загружаются; политика FSM содержит 11 состояний; вероятности переходов суммируются в \(1{,}0\) \\
\texttt{test\_preprocess.py} & shape выходов; отсутствие NaN/Inf; неизменность статистик при повторном fit; save/load roundtrip; обработка неизвестной категории; обрезка выбросов; перенос меток \\
\texttt{test\_windowing.py} & корректные shapes; стратегии \emph{any}/\emph{majority}/\emph{last}; обработка коротких последовательностей; валидация входов \\
\texttt{test\_splits.py}    & сохранение временного порядка; воспроизводимость при random; непересечение split-ов \\
\texttt{test\_models\_smoke.py} & все 10 DL-моделей загружаются из YAML и выполняют forward; classical baselines fit/predict\_proba; save/load CNN-LSTM; CNN-LSTM помечен как \emph{proposed} \\
\texttt{test\_metrics.py}  & Precision/Recall/F1/MCC, PR-AUC/ROC-AUC; threshold для целевого FPR; ECE; калибровка снижает overconfidence; PSI/KL/MMD на одинаковых и сдвинутых распределениях; bootstrap CI содержит выборочную метрику; FocalLoss конечен и положителен \\
\texttt{test\_attacker.py} & генерация всех 7 шаблонов; FSM-агент эмиссия тиков; drift-инжектор сдвигает признаки и маскирует сигнатуры; runtime даёт DataFrame; воспроизводимость по seed \\
\texttt{test\_alerting.py} & уровни severity ladder; пропуск под порогом; dedup в окне; обновление history до history\_size \\
\bottomrule
\end{tabularx}
\end{table}

\section{Соответствие постановке задачи и решениям Этапов~1--2}

В \tabref{tab:traceability} зафиксировано соответствие реализованных
артефактов требованиям постановки и проектным решениям.

\begin{table}[H]
\centering
\caption{Соответствие реализации проектным решениям}
\label{tab:traceability}
\renewcommand{\arraystretch}{1.2}
\small
\begin{tabularx}{\linewidth}{@{}lXl@{}}
\toprule
\textbf{Требование (Этап~1, §4)} & \textbf{Артефакт реализации} & \textbf{Файл / тест} \\
\midrule
FR-1 (схема UNSW-NB15) & \texttt{UNSWConfig}, \texttt{UNSWFeatureSchema} & \texttt{schema.py}, \texttt{test\_configs.py} \\
FR-2 (препроцессинг) & \texttt{Preprocessor}, \texttt{build\_windows} & \texttt{preprocess.py}, \texttt{test\_preprocess.py} \\
FR-3 (\(\geq 14\) моделей) & регистр \texttt{models}; 10 DL + 5 classical & \texttt{models/}, \texttt{test\_models\_smoke.py} \\
FR-4 (порог + калибровка) & \texttt{find\_threshold\_for\_target\_fpr}, \texttt{TemperatureScaler} & \texttt{eval/}, \texttt{test\_metrics.py} \\
FR-5 (алерты) & \texttt{AlertFormer}, severity ladder & \texttt{inference/alerting.py}, \texttt{test\_alerting.py} \\
FR-6 (REST) & FastAPI \texttt{POST /score} & \texttt{inference/service.py} \\
FR-7 (UI) & Streamlit dashboard & \texttt{ui/streamlit\_app.py} \\
FR-8 (агент) & \texttt{FSMAgent + DriftInjector + AttackerRuntime} & \texttt{attacker/}, \texttt{test\_attacker.py} \\
FR-9 (drift-monitor) & \texttt{drift\_report}, PSI/KL/MMD & \texttt{eval/drift.py}, \texttt{test\_metrics.py} \\
FR-10 (seed журнал) & \texttt{set\_seed}, JSON-логи в \texttt{experiments/} & \texttt{utils/seed.py}, \texttt{test\_utils.py} \\
\bottomrule
\end{tabularx}
\end{table}

\section{Промежуточные выводы}

\begin{enumerate}
    \item Программная реализация прототипа выполнена в полном объёме и
    соответствует проектным решениям, зафиксированным в Этапах~1 и 2.
    \item Поддерживается весь набор моделей-кандидатов из протокола
    Эксперимента~1: CNN-LSTM (\emph{proposed}), 9 альтернативных DL-моделей
    и 5 классических baselines, что обеспечивает условия для сравнения в
    равных условиях.
    \item Реализован ИИ-злоумышленник в архитектуре \emph{templates +
    FSM + drift}, поддерживающий offline и online режимы и обеспечивающий
    воспроизводимую ground truth.
    \item Realtime-стенд (FastAPI + Streamlit) пригоден для проведения
    Экспериментов~4--6 Этапа~4 и для демонстрации работы прототипа в
    рамках защиты ВКР.
    \item Покрытие тестами включает 9 тестовых модулей; критические
    инварианты (отсутствие утечки в препроцессоре, корректность шейпов
    окон, суммирование переходов FSM в \(1{,}0\), снижение PSI на
    идентичных распределениях, dedup алертов) проверены формально.
\end{enumerate}

% ----------- Список использованных источников -----------
\printbibliography[title={СПИСОК ИСПОЛЬЗОВАННЫХ ИСТОЧНИКОВ}]

\label{LastPage}

\end{document}
